% C.3bj finalized after complete owner, root and both peer proof/source audits.
\paragraph{C.3bj (Proved): integral higher kernels in every positive characteristic.}
Let $k$ be an algebraically closed field of characteristic $p>0$ and
let $N\ge5$. There exist an excellent normal strictly henselian local
$k$-domain $R$ of dimension $N$, with residue field $k$, and
$0\ne f\in\mathfrak m_R$, such that the puncture $X$ and the principal
punctured divisor $Y=X\cap V(f)$ are regular and
\[
 \ker\bigl(\mathrm{CH}^2(X)_{\rm alg}
       \longrightarrow\mathrm{CH}^2(Y)_{\rm alg}\bigr)
\text{ contains a class of exact order }p.
 \tag{C.3bj.1}
\]
Here ordinary algebraic equivalence means actual families on fixed
products with smooth proper curves. The data can be chosen over
$\overline{\mathbb F}_p$ before extension of the algebraically closed
field. This statement concerns integral coefficients, including $p=2$;
it asserts no nonzero rational kernel. C.3bi retains its independent
specified characteristic-five construction and its different proof.

\begin{proof}
\textbf{1. The projective data and the nonzero divisor cup.}
Choose a smooth ordinary plane quartic $C=(F=0)$ and a plane quartic
$H$ for which the Serre-dual normal-deformation form
\[
 Q_H(s,t)=\langle\delta(H),st\rangle,
 \qquad s,t\in H^0(C,K_C),\qquad K_C=\mathcal O_C(1),
 \tag{C.3bj.2}
\]
is nondegenerate. C.3bg proves existence in every characteristic,
including two, and gives a terminating finite-field enumeration
specifying such a pair. Recall the two essential computations:
$z_0z_1z_2(z_0+z_1+z_2)$ has Hasse--Witt matrix $I_3$ in every
characteristic, so ordinarity is a nonempty open condition; and the
Euler and normal sequences identify the surjective normal-deformation
map with the dual of
$\operatorname{Sym}^2H^0(K_C)\simeq H^0(K_C^2)$.
Every symmetric bilinear form therefore occurs, without division by two.
Smoothness is another nonempty open condition. We choose the first
finite-field pair in the enumeration of C.3bg satisfying these tests.

Let $W$ be a smooth cubic surface, put $A=\mathcal O_W(1)$, and
choose a line $D\subset W$. These data may also be specified over a
finite field. For $p\ne3$ take the Fermat cubic with
$D=(w_0+w_1=w_2+w_3=0)$. In characteristic three take
\[
 W:\ u s^2+v t^2+u^2t+v^2s=0,\qquad D=(u=v=0).
 \tag{C.3bj.3}
\]
For this last surface, a common zero of the four derivatives with a
zero coordinate forces all coordinates to vanish. Otherwise set $u=1$;
the derivative equations give $s=v^2$, $t=v^{-1}$ and $v^5=1$,
whereas the equation of the surface becomes $4v^4=v^4\ne0$.
Thus it is smooth. In all cases $A=-K_W$, $D^2=-1$ and $A\cdot D=1$.

The twisted divisor cup map
\[
 \mu_D:H^0(W,A)\longrightarrow H^1(W,A\otimes\Omega^1_{W/k}),
 \qquad a\longmapsto a\,d\log\mathcal O_W(D),
 \tag{C.3bj.4}
\]
has rank two in every characteristic. Here is the characteristic-free
normalization check. Write $D=(u=v=0)$ and $F_W=uQ_0+vQ_1$.
Smoothness along $D$ says that the two binary quadratics
$Q_0|_D,Q_1|_D$ have no common zero. The Euler and conormal sequences
identify the target of (C.3bj.4) with $H^0(W,A)^\vee$.
The twisted normal sequence of the line is
\[
 0\longrightarrow\mathcal O_D(-2)\longrightarrow\mathcal O_D^2
       \xrightarrow{(Q_0,Q_1)}\mathcal O_D(2)\longrightarrow0.
\]
The boundary to $H^2(W,\mathcal O_W(-A))$ is an isomorphism:
$H^1$ and $H^2$ of $\mathcal O_W(D-A)$ vanish by Riemann--Roch
and Serre duality, as calculated in C.3bi. Lifting the local defining
equations of the line identifies this boundary with contraction
against its $d\log$ class. Consequently the dual bilinear form is,
up to a single nonzero common scalar,
\[
 (a,b)\longmapsto\lambda_D((ab)|_D),
 \qquad
 \ker\lambda_D=\langle Q_0|_D,Q_1|_D\rangle.
 \tag{C.3bj.5}
\]
This two-variable symmetric form is nonsingular: a rank-one symmetric
form over an algebraically closed field is a scalar evaluation form,
also in characteristic two. If (C.3bj.5) were degenerate, both
quadratics would vanish at the corresponding point of $D$, which
contradicts their having no common zero. Equivalently, in our specified
models the two quadratics are nonzero scalar multiples of $s^2,t^2$,
and the form has matrix $\left(\begin{smallmatrix}0&1\\1&0\end{smallmatrix}\right)$.
No factor of two occurs.

Put $m=N-4\ge1$, $S=W\times C$ and $B=A\boxtimes K_C$.
Choose $a_0,a_1\in H^0(S,B)$ with smooth codimension-two common
zero scheme, and write
\[
 x_0a_0+x_1a_1=\sum_{j=0}^2g_j z_j,
 \qquad g_j\in H^0(\mathbb P^1,\mathcal O(1))\otimes H^0(W,A).
\]
Require that $g_0,g_1,g_2$ be independent and that their span have
nonzero image under $1\otimes\mu_D$. These are simultaneous nonempty
open conditions on the affine space of ordered pairs. Smoothness is
Bertini for the very ample $B$, successively on its smooth first
section \citep[Lemma 33.47.3, Tag 0FD6]{D_StacksBertiniUniformHigher}.
The two rank conditions are nonempty because, using the coordinates
$s,t$ on $D$, the three sections $x_0s,x_0t,x_1s$ give a witness.
Irreducibility of the parameter space makes their intersection nonempty.

If $m\ge2$, choose $a_2$ whose zero divisor on the preceding smooth
base curve is reduced. If $m\ge3$, choose $a_3$ nonzero at every
point of the finite scheme $(a_0=a_1=a_2=0)$; put $a_i=0$ for
$i\ge4$. Define
\[
 V=\mathbb P^m\times W\times C,\qquad
 L=\mathcal O_{\mathbb P^m}(1)\boxtimes A\boxtimes K_C,
 \qquad\ell=\sum_{i=0}^m x_i a_i.
 \tag{C.3bj.6}
\]
All choices can be made over a finite extension of the coefficient
field of $C,W$: the stated nonempty open sets have closed points.
They can be specified by enumerating finite extensions and coefficient
tuples and selecting the first tuple passing the displayed Jacobian,
rank and avoidance tests. This is a terminating specification, not a
claim of a uniform bound on search time.
The incidence divisor $T=(\ell=0)\subset V$ is smooth. Off the common
base locus, its projective-space derivative is nonzero. At a base
point the differentials of the $m+1$ sections are independent for
$m=1,2$, so their combination with a nonzero vector $(x_i)$ is nonzero.
For $m\ge3$ there is no common base point. The original three
coefficients remain independent after adjoining the extra $x_i$,
and their divisor cup remains nonzero by restriction to $\mathbb P^1$.

\textbf{2. The actual local ring and the split principal cut.}
In the total space of
$\mathcal O_{\mathbb P^m\times W\times\mathbb P^2}(-1,-1,-1)$
take the equation
\[
 Z_0:\quad F+\tau\ell H=0.
 \tag{C.3bj.7}
\]
Contract by all the Segre coordinates $\tau x_iw_jz_h$, take the
reduced affine image, normalize it, localize at the unique vertex,
and strictly henselize. This defines $R$ from the preceding specified
data; take $f=\tau\ell$, the corresponding linear combination of
those Segre coordinates. The graph and proper contraction proof in
C.3bf--C.3bg applies with the additional smooth factor $W$:
$F$ does not divide $\ell H$, the equation is integral, and $dF$
is nonzero along $C$. The exceptional fibre is the connected smooth
$V$, normalization has one point above the vertex, and the contraction
is an isomorphism off that fibre. Its singular locus away from a regular
neighbourhood of $V$ has proper image avoiding the vertex and disappears
after this localization. Thus the induced projective map
$\mathcal Z\to\operatorname{Spec}R$ has regular source $\mathcal Z$,
is an isomorphism off $V$, and satisfies
\[
 \dim R=m+4=N,\qquad
 \mathfrak m_R\mathcal O_{\mathcal Z}=I_V,\qquad I_V/I_V^2=L.
 \tag{C.3bj.8}
\]
The maximal-ideal identity follows because the Segre generators generate
the fibre ideal on every projective coordinate chart. The projection
$\mathcal Z\to W$ is retained, so $\mathcal D:=\mathcal O_W(D)$
pulls back to an actual line bundle on $\mathcal Z$.
On the strict transform of $f=0$, equation (C.3bj.7) becomes exactly
$F=0$; hence this transform is the split line bundle with zero section
$T$ and conormal $L|_T$. Consequently both $X=\mathcal Z-V$ and its
principal cut $Y$ are regular.

\textbf{3. An ordinary Picard endpoint with nonzero divisor cup.}
The cubic and curve cohomology and K\"unneth give
\[
 \begin{split}
 H^1(V,L)&=H^0(\mathbb P^m,\mathcal O(1))\otimes H^0(W,A)
                       \otimes H^1(C,K_C),\\
 H^2(V,L)&=0,\qquad
 H^1(V,L^n)=H^2(V,L^n)=0\quad(n\ge2),\\
 H^3(V,L^n)&=0\quad(n\ge0).
 \end{split}
 \tag{C.3bj.9}
\]
Let $E=V_2$ be the first infinitesimal neighbourhood and
$J=\operatorname{Pic}^0(C)$. Its Picard extension is
\[
 0\longrightarrow U_0:=H^1(V,L)\longrightarrow G
                      \longrightarrow J\longrightarrow0.
 \tag{C.3bj.10}
\]
The first-neighbourhood/Atiyah calculation of C.3bg, with $W$ carried
along, identifies its extension-class map with
\[
 s\longmapsto\sum_jQ_H(z_j,s)g_j.
\]
Its image is the rank-three space
$U=\ell H^1(V,\mathcal O_V)=\langle g_0,g_1,g_2\rangle$,
where now $g_j$ denotes the full coefficient in all the $x_i$.
This is the characteristic-independent transition-unit formula of
\citet[Theorem 2.6 and formula (2.7.2), pp.~166--167]{D_Illusie1972Torsors}.
The cubic divisor cup is nonzero on $U$ by Step 1 and the direct
K\"unneth summand
\[
 H^0(\mathbb P^m,\mathcal O(1))\otimes
 H^1(W,A\otimes\Omega^1_{W/k})\otimes H^1(C,K_C)
 \subset H^2(V,L\otimes\Omega^1_{V/k}).
 \tag{C.3bj.11}
\]

Ordinarity gives the etale Verschiebung $v:J^{(p)}\to J$, with
$K=\ker v\simeq(\mathbb Z/p)^3$ and $v^*H^1(J,\mathcal O_J)=0$.
The integral extension classification
$\operatorname{Ext}^1(J,\mathbb G_a)=H^1(J,\mathcal O_J)$
\citep[Section 3.2.3, equation (3.8), p.~23]{D_BrionIsogenyII}
gives a homomorphism $b:J^{(p)}\to G$ over $v$.
The Hom--Ext sequence shows that $b(K)$ spans $U$ and that $b|_K$
is injective. Thus one can choose $\eta\in K$ such that
\[
 a=b(\eta),\qquad
 a\smile d\log\mathcal D\ne0\quad
       \text{in }H^2(V,L\otimes\Omega^1_{V/k}).
 \tag{C.3bj.12}
\]

The target calculation needs no identification of its entire Picard
variety. The divisor sequence and (C.3bj.9), including
$H^2(V,\mathcal O_V)=0$, give
\[
 H^1(T,L|_T)=U_0/U,\qquad
 H^1(T,L^n|_T)=H^2(T,L^n|_T)=0\quad(n\ge2).
 \tag{C.3bj.13}
\]
The first target neighbourhood is split. Its Picard functor is therefore
the product of that of $T$ and the vector group $U_0/U$: pullback
along the retraction splits the unit sequence. The first coordinate
of the restriction of $b$ factors through $v$ and the pullback from
$C$ to $T$. Its vector coordinate is zero, since a morphism from the
proper connected $J^{(p)}$ to a vector group is constant and the
identity maps to zero. Thus $a=b(\eta)$ has zero target first jet;
(C.3bj.13) makes every higher endpoint lift unique.

Choose a smooth proper connected curve $P\subset J^{(p)}$ through
$0$ and $\eta$. Pull the normalized Poincare bundle to $V\times P$
through $v|_P$. Its first lifting obstruction is the mixed Atiyah
class, whose $H^1(J,\mathcal O_J)$ factor is killed by $v^*$.
It therefore lifts to $E\times P$. Its classifying map differs from
$b|_P$ by a constant map to $U_0$; tensor by the inverse fibre at zero
to remove that constant. For $n\ge2$, (C.3bj.9) and
$H^2(P,\mathcal O_P)=0$ give $H^2(V\times P,L^n)=0$.
The bundle lifts compatibly through all exceptional thickenings.
Proper formal existence on $\mathcal Z_{\widehat R}\times P$,
finite-data Popescu descent, and smooth henselian lifting modulo
$\mathfrak m_R^2$ produce an actual line bundle on the fixed product
$\mathcal Z\times P$ with precisely this first jet. Include its inverse
and the first-jet isomorphism among the finite descent data, as in
C.3bf and C.3bg. Normalize its fibre at zero to be trivial on
$\mathcal Z$, and let $\mathcal A$ be the fibre at $\eta$.
Then $\mathcal A|_V=\mathcal O_V$ and its relative first jet is $a$.
Higher endpoint lifts are unique by (C.3bj.9); the same is true on
the target by (C.3bj.13). Proper formal existence and Elkik's injective
punctured Picard comparison \citep[unnumbered corollary, p.~573]{D_Elkik1973}
therefore prove
\[
 \mathcal A|_X\in\operatorname{Pic}(X)_{\rm alg},\qquad
 (\mathcal A|_X)^p\simeq\mathcal O_X,\qquad
 \mathcal A|_Y\simeq\mathcal O_Y.
 \tag{C.3bj.14}
\]
This uses actual curve families, not a family on a merely formal or
henselized product.

\textbf{4. The precise divisible group before the relative detector.}
For a smooth projective variety $Q$ put
$G^{a,b}(Q)=H^a(Q,\mathcal K_b^M)$, using the ordinary Milnor
symbol sheaf. Its integral Gersten complex identifies it with Rost's
groups with Milnor coefficients
\citep[Theorems 7.1 and 7.5, author pp.~20--22]{cyKerz2009Milnor}.
Rost localization, affine homotopy and pullback give a surjection
\[
 G^{1,2}(Q)\oplus k^*\twoheadrightarrow
                 G^{1,2}(\mathbb P^d\times Q)\quad(d\ge1).
\]
Indeed remove $\mathbb P^{d-1}\times Q$: the open group is
$G^{1,2}(Q)$, lifted by projection, and the remaining image comes
from $G^{0,1}(\mathbb P^{d-1}\times Q)=k^*$.
For a blow-up $\rho:\widetilde Q\to Q$ along a smooth connected
centre, the class $c-\rho^*\rho_*c$ vanishes away from the exceptional
divisor and hence comes from its $G^{0,1}=k^*$. The exact integral
inputs are \citet[Proposition 4.1(3), Section 5, Corollary 6.5,
Proposition 8.6, and Section 12, pp.~351, 355--360, 370, 382--383]{cyRost1996}.
Every smooth cubic surface over an algebraically closed field is a
blow-up of $\mathbb P^2$ in six points
\citep[Chapter 4, Proposition 2.4, p.~178]{D_HuybrechtsCubics}.
After product with $C$, these operations prove that $H^1(V,\mathcal K_2^M)$
is a quotient of
\[
 H^1(C,\mathcal K_2^M)\oplus(k^*)^8.
\tag{C.3bj.15}
\]
The curve group is the cokernel of
$K_2^M(k(C))\to\bigoplus_{c\in C(k)}k^*$.
Algebraic closedness makes $k^*$ divisible; direct sums and quotients
of divisible groups are divisible. Thus $H^1(V,\mathcal K_2^M)$ is
divisible. This is a sheaf-cohomology assertion, not an inference
from the ordinary group $K_1(V)$.

\textbf{5. A canonical integral relative Milnor class, including $p=2$.}
Regard $E$ and $V$ as having the same underlying Zariski space, and put
\[
 \mathcal J=\ker(\mathcal K_2^M(E)\to\mathcal K_2^M(V)).
\]
Units lift across a square-zero quotient. Every downstairs Steinberg
relation lifts as $\{u,1-u\}$, with both entries still units.
Consequently the Milnor map is onto and its kernel is generated
locally by symbols having an entry $1+i$, $i\in I=L$.
Since $(1+i)^p=1$, multilinearity gives $p\mathcal J=0$.
This argument includes characteristic two and keeps symbols with two
infinitesimal entries; none is silently declared zero.

Smoothness of $V$ gives local algebra splittings of its square-zero
thickening \citep[Definition 10.138.1 and Proposition 10.138.13]{frStacksFormalSmoothRelativeMilnor}. The conormal sequence is therefore canonically short exact
\[
 0\longrightarrow I\xrightarrow{d}\Omega^1_{E/k}|_V
             \longrightarrow\Omega^1_{V/k}\longrightarrow0.
 \tag{C.3bj.16}
\]
The usual conormal sequence is \citet[Lemma 10.131.9, Tag 00RU]{D_StacksConormalUniformHigher};
its left injectivity here follows directly on a splitting
$\mathcal O_V\oplus I$, without a regular-immersion assumption on $V\subset E$.
Since $I$ has rank one, the kernel on second exterior powers is
canonically $I\otimes\Omega^1_{V/k}$, by $i\otimes\omega\mapsto di\wedge\widetilde\omega$.
The ordinary $d\log$ of Milnor symbols restricts for $\mathcal J$ into
this kernel, and hence defines a canonical map
\[
 \tau:\mathcal J\longrightarrow I\otimes\Omega^1_{V/k},\qquad
 \{1+i,u\}\longmapsto i\otimes d\log\bar u.
 \tag{C.3bj.17}
\]
It is independent of a local splitting and of a frame for $I$.
Symbols with both entries infinitesimal have zero image because the
wedge of two sections of the rank-one conormal submodule is zero.
The Steinberg relation is respected by $du\wedge d(1-u)=0$,
also in characteristic two.

The long exact sequence of the exact Milnor sheaf sequence gives
\[
 H^1(V,\mathcal K_2^M)\longrightarrow H^2(V,\mathcal J)
       \longrightarrow H^2(E,\mathcal K_2^M).
\]
The first map is zero: its source is divisible and its target is
killed by $p$. Hence the second map is injective. With local
trivializations of $\mathcal A|_E$ respecting its trivialization on
$V$, write its transitions $1+a_{ij}$ and those of $\mathcal D$ as
$g_{ij}$. The cup cocycle $\{1+a_{ij},g_{jk}\}$ defines a class
$\xi\in H^2(V,\mathcal J)$ with
\[
 H^2(\tau)(\xi)=a\smile d\log\mathcal D\ne0.
 \tag{C.3bj.18}
\]
Its image in $H^2(E,\mathcal K_2^M)$ is exactly the absolute cup
product of the two actual line-bundle classes, and is nonzero.
No Milnor--Quillen comparison on the nonreduced $E$, no rational
Chern character, and no Adams-weight separation enter this argument.

\textbf{6. Removing the whole exceptional image and concluding.}
Let $\beta=c_1(\mathcal A)c_1(\mathcal D)\in\mathrm{CH}^2(\mathcal Z)$.
Its restriction to $X$ is ordinary algebraically trivial, by the actual
curve family of Step 3 multiplied by the fixed actual divisor $D$.
It is killed by $p$ and restricts to zero on $Y$ by (C.3bj.14).
Suppose it vanished on $X$. Localization would give
$\beta=i_*\gamma$ with $\gamma\in\mathrm{CH}^1(V)=\operatorname{Pic}(V)$.
Restricting to $V$, where $\mathcal A$ is trivial, and using the
self-intersection formula gives $c_1(L)\gamma=0$.
We now check injectivity of this multiplication integrally, retaining
all possible torsion in $\operatorname{Pic}(C)$.
Write $h=c_1(\mathcal O_{\mathbb P^m}(1))$, $c=c_1(B)$ and
$\gamma=\gamma_0+bh$ with $\gamma_0\in\operatorname{Pic}(S)$,
$b\in\mathbb Z$. If $m\ge2$, the coefficient of $h^2$ in
$(h+c)(\gamma_0+bh)$ is $b$, and its coefficient of $h$ is
$\gamma_0+bc$; thus $b=0$ and $\gamma_0=0$.
If $m=1$, the projective-bundle coefficients instead give
$\gamma_0=-bc$ and $bc^2=0$. Intersect with $c$ and take degree on
$S$: $\deg(c^3)=3\deg(A^2)\deg(K_C)=36$ as an integer, so $b=0$
and again $\gamma_0=0$. This integer degree argument remains valid
when $p$ divides $36$; Chow coefficients have not been reduced modulo $p$.
Therefore $\beta=0$ on $\mathcal Z$.

We apply the integral Milnor--Bloch formula and its product
compatibility on an actual smooth finite-type model. Descend both
line bundles and the hypothesized relation $\beta=0$, including the
finitely many cycles and rational functions witnessing it, to one
pointed etale neighbourhood of the vertex. This is possible because
all of these data are finite and the strict henselization is the
filtered limit of those neighbourhoods. Shrink the affine base around
its distinguished point so that the corresponding proper model of
$\mathcal Z$ is regular: its nonregular locus has proper closed image
avoiding that point. This model is smooth over the perfect field $k$.
The etale local map has the same residue field $k$ and identifies the
quotients by the square of the maximal ideal. Thus its exceptional
first jet, the two restricted bundles, and the nonzero class computed
on $E$ are exactly the ones in Step 5; they have not been replaced by
an unspecified formal model. Codimension and the descended rational
relation are preserved at a sufficiently late etale stage.

On this smooth finite-type model the integral Milnor--Bloch formula
and its product compatibility identify $\beta$ with the cup of the
two line-bundle classes
\citep[Theorems 7.1 and 7.5, author pp.~21--22]{cyKerz2009Milnor}
\citep[Section 3, diagram (22), pp.~557--558]{cyBraunling2013}.
If $\beta=0$, that cup is zero. Its natural restriction to the actual
nonreduced $E$ would therefore be zero, contradicting Step 5.
This finite-model argument requires neither a Bloch formula on $E$
nor an unstated continuity theorem for second Milnor-sheaf cohomology.
Thus $\beta|_X\ne0$, and (C.3bj.14) makes its order exactly $p$.
All the constructions and nonvanishing calculations remain valid
under extension to any algebraically closed field of characteristic
$p$, proving (C.3bj.1).
\end{proof}


\noindent\emph{Integral codimension-two injectivity over the full class
of coefficient fields.} For each prime $p$, universal injectivity of
ordinary integral codimension-two restriction, over all algebraically
closed coefficient fields of characteristic $p$ in the stated local
geometric scope, holds exactly when $N\le2$. In those dimensions the
regular puncture has no codimension-two points. For $N=3$ and $N=4$,
C.3aw and C.3ay with $r=2$ respectively give counterexamples over their
specified transcendental coefficient fields. For every $N\ge5$,
C.3bj gives a counterexample already over $\overline{\mathbb F}_p$
and over every algebraically closed extension. This corollary does
not assert that the lower-dimensional counterexamples from C.3aw
and C.3ay occur over every fixed algebraically closed coefficient field.
